\documentclass[11pt]{article}
\usepackage[T1]{fontenc}
\usepackage[utf8]{inputenc}
\usepackage{lmodern,amsmath,amssymb,amsthm,booktabs,array,microtype}
\usepackage[margin=25mm]{geometry}
\usepackage[hidelinks]{hyperref}
\usepackage{xurl}
\newtheorem{theorem}{Theorem}
\newtheorem{lemma}[theorem]{Lemma}
\newtheorem{proposition}[theorem]{Proposition}
\allowdisplaybreaks
\setlength{\emergencystretch}{2em}
\title{Verification supplement for the conditional Subset Sum refinement}
\author{Yanni Bills}
\date{21 October 2026}
\begin{document}
\maketitle
\begin{abstract}
We separate the obligations behind the published conditional exponent
$0.4876723$ and reconstruct their concrete mathematical interfaces.
The supplement gives full written proofs of the filter-3 boundary
transfer, analytic score bounds and integration error method, effective
finite-circle approximation, and the isolation/preliminary-search
case split. All 430 original continuous moment inequalities have now
been validated and independently replayed using explicit inner/outer
error bounds and ball arithmetic outside Lean. A separately reviewed
Lean development establishes concrete boundary and regularity results,
exact data links, real-logarithm and entropy bounds, all-row conditional
transfers and unique region assignment. The remaining written links and
end-to-end algorithmic obligations are distinguished from these results. The published release is preserved unchanged;
no stronger unconditional running-time claim is made here.
\end{abstract}

\section{Baseline, obligations and evidence levels}
The baseline is public release \texttt{v1.0.0}, commit
\texttt{5380ae2c6411f70769df6076ca82a54d43ff09c3} of
\url{https://github.com/yannouli/subset-sum-conditional-refinement}.
All additions live under \texttt{verification/}; the released manuscript,
tuple, code, pins and historical evidence remain unchanged.
The baseline assumes original algorithmic claims and 430 continuous
moment inequalities. Publication, a Claimed listing, old Python checks
and internal review are not substitutes for proofs of those assumptions.

The four groups are: (A) our changed boundary/transfer argument;
(B) exact fixed data, real logarithm bounds and scalar allowances;
(C) actual analytic scores, continuous moments and finite-circle limits;
and (D) algorithmic correctness, precision and every-history word cost.
The machine-readable obligation record is
\texttt{theory/obligations.json}; intended statements are stated before
implementation in \texttt{theory/intended-statements.md}.
The accepted numerical evidence is described in
Section~\ref{sec:validated-results}. The separately accepted concrete
Lean checkpoint is mapped precisely in Section~\ref{sec:formal-results};
its limitations and the unfinished stronger replay are explicit.

\paragraph{Before and after the written reconstruction.}
The score bound and finite-circle approximation need not be retained
as unexamined inherited assertions: Section~\ref{sec:analytic} and
Section~\ref{sec:finite} prove them for the concrete definitions.
Section~\ref{sec:algorithmic} similarly reconstructs the isolation and
LD/preliminary-search implications. The moment inequalities are now
supported by validated numerical evidence with explicit computational
trust; the algorithmic dependencies in Section~\ref{sec:remaining}
remain separate. Validation of every continuous moment does not remove
those remaining algorithmic obligations. No review badge is inferred.

\section{The actual analytic functions}\label{sec:analytic}
All integrals in this section use the probability average
\[
 \langle h\rangle=\frac1{2\pi}\int_0^{2\pi}h(t)\,dt.
\]
In particular this is not an integral over a region selected by a pattern.
For each of the 43 rational coefficient records of Kapila's Appendix~A,
put
\[
 r_v(x)=\sum_{j=1}^8r_{v,j}\cos(jx),\quad
 s_v(x)=\sum_{j=1}^8s_{v,j}\cos(jx),\quad
 F_v=\exp((\ln2)r_v),\quad G_v=\exp((\ln2)s_v).
\]
Suppress $v$ temporarily. The exponent $\kappa=\kappa_v$ is the exact
source rational; $1\le\kappa\le2$. Define
\begin{align*}
 Z_f&=\langle F\rangle,& Z_g&=\langle G\rangle,&
 Z_{f,\kappa}&=\langle F^\kappa\rangle,&
 Z_{g,\kappa}&=\langle G^\kappa\rangle,\\
 f&=F/Z_f,&g&=G/Z_g,& f_*&=F^\kappa/Z_{f,\kappa},&
 g_*&=G^\kappa/Z_{g,\kappa}.
\end{align*}
\begin{align*}
 \mathcal P(x)&=\langle F(x+t)G(t)\rangle_t,\\
 \mathcal Y(x)&=\langle F(x+t)G(t)(r(x+t)+s(t))\rangle_t,\\
 Q(x)&=\mathcal P(x)/(Z_fZ_g),\\
 K_x(t)&=f(x+t)g(t)/Q(x).
\end{align*}
The intended scores, following the definitions rather than their names,
are
\begin{align*}
 D(x)&=\langle K_x\log_2K_x\rangle,&
 A_s(x)&=\langle K_x(t)\log_2f_*(x+t)\rangle_t,\\
 E_s(x)&=\langle K_x\log_2g_*\rangle,&
 u(x)&=A_s(x)+E_s(x)-D(x),\quad
 w(x)=2D(x)-A_s(x)-E_s(x).
\end{align*}
The subscripts distinguish scores $A_s,E_s$ from a budget ratio or an
exponent. All logarithms in scores are real base-two logarithms.

\begin{lemma}[Concrete existence and score bound]\label{lem:exist}
For these definitions all denominators are strictly positive, the four
density normalizations and $K_x$ have average one, and all real scores
are continuous and bounded. Every affine-certificate moment integrand
is continuous and integrable. If each coefficient has absolute value
at most two, then $|w(x)|\le256$ for every $x$.
\end{lemma}
\begin{proof}
The coefficient assumption gives $|r|,|s|\le16$. Hence
$2^{-16}\le F,G,Z_f,Z_g\le2^{16}$ and
$2^{-32}\le\mathcal P\le2^{32}$. The tilted quantities satisfy the
corresponding bounds with $16\kappa$ in place of 16.
All assertions of positivity follow before taking any logarithm or
quotient. Periodicity and substitution show that the averages of the
normalized functions are one. The average of $K_x$ is one by the
definition of $Q$. Integration of continuous functions depending
continuously on $x$ on a compact rectangle preserves continuity;
the strictly positive denominator bounds preserve it through division
and logarithms. This proves integrability as well as continuity, and
prevents a totalized logarithm or integral from masking a missing hypothesis.

Put
\[
 L=\log_2\mathcal P-\log_2(Z_fZ_g),\quad
 B=\mathcal Y/\mathcal P-\log_2(Z_fZ_g),\quad
 H=\log_2(Z_{f,\kappa}Z_{g,\kappa})-\kappa\log_2(Z_fZ_g).
\]
Here $B$ is the $K_x$-average of $r(x+t)+s(t)-\log_2(Z_fZ_g)$.
Substitution gives $D=B-L$ and $A_s+E_s=\kappa B-H$, so
\begin{equation}\label{eq:score-identities}
 u=L+(\kappa-1)B-H,\qquad w=(2-\kappa)B-2L+H.
\end{equation}
The quotient $\mathcal Y/\mathcal P$ is an average of $r+s$ with
positive weights summing to one, hence has absolute value at most 32.
Also $|\log_2(Z_fZ_g)|\le32$ and $|\log_2\mathcal P|\le32$.
Thus $|B|,|L|\le64$. The tilted normalizer bounds give $|H|\le64\kappa$.
Since $0\le2-\kappa\le1$,
\[
 |w|\le64(2-\kappa)+128+64\kappa=256.
\]
This derives the score bound through the actual definitions, rather
than stopping at a bound on trigonometric coefficients.
\end{proof}

\begin{lemma}[The required entropy order]\label{lem:entropy-order}
The actual scores satisfy $D\ge0$, $D-A_s\ge0$, $D-E_s\ge0$,
and hence $w\ge0$. The same is true for their exact finite-circle
versions.
\end{lemma}
\begin{proof}
If $p,q$ are strictly positive normalized densities, the elementary
inequality $\ln t\le t-1$ gives
\[
 \langle p\ln(q/p)\rangle\le\langle q-p\rangle=0.
\]
Apply this with $p=K_x$ and $q=1,f_*(x+\cdot),g_*$ respectively,
and divide by $\ln2>0$. Each $q$ is normalized; for the shifted
one this follows from periodic translation invariance. The same
argument uses only normalized sums on the finite circle, where grid
translations permute the summands. Finally
$w=(D-A_s)+(D-E_s)\ge0$.
\end{proof}

\paragraph{Independent normalization check.}
Eliminating all untilted normalizers from~\eqref{eq:score-identities}
gives the exact equivalent formulas
\begin{align}
 u&=\log_2\mathcal P+(\kappa-1)\mathcal Y/\mathcal P
       -\log_2(Z_{f,\kappa}Z_{g,\kappa}),\label{eq:cancel-u}\\
 w&=(2-\kappa)\mathcal Y/\mathcal P-2\log_2\mathcal P
       +\log_2(Z_{f,\kappa}Z_{g,\kappa}).\label{eq:cancel-w}
\end{align}
The tilted normalizers do not cancel. These identities are useful for
an independent implementation check, not alternative definitions that
avoid the density correspondence.

\section{A validated integration contract}\label{sec:quadrature}
This section proves a sufficient error bound. It does not assert that
any numerical inequality has passed until the specified enclosures
have actually been produced and checked. All source decimals are
exact rationals, including coefficients and multipliers.

\begin{lemma}[Periodic trapezoid bound]\label{lem:trap}
Suppose a $2\pi$-periodic function $h$ is holomorphic on a neighborhood
of the closed strip $|\operatorname{Im}z|\le\sigma$, where $\sigma>0$,
and $|h|\le B_0$ there. For $N\ge1$,
\[
 \left|\frac1N\sum_{j=0}^{N-1}h(2\pi j/N)-\langle h\rangle\right|
 \le \mathcal E_N(B_0,\sigma):=\frac{2B_0}{e^{N\sigma}-1}.
\]
\end{lemma}
\begin{proof}
For a nonzero Fourier index $k$, shift the integral for its coefficient
to $\operatorname{Im}z=-\sigma\operatorname{sgn}k$. Periodicity cancels
the vertical contour sides and gives
$|\widehat h(k)|\le B_0e^{-|k|\sigma}$. Thus the Fourier series is
absolutely convergent on the real axis. Averaging its terms at the
$N$ nodes kills every index not divisible by $N$. The error is at most
$2B_0\sum_{j\ge1}e^{-jN\sigma}$, the claimed geometric sum.
\end{proof}

\paragraph{Inner integrals and arithmetic.}
Define $R(t)=\sum_j|r_j|\cosh(jt)$ and
$S(t)=\sum_j|s_j|\cosh(jt)$. For an inner strip width $\sigma>0$,
$|r(z)|\le R(\sigma)$ and $|s(z)|\le S(\sigma)$.
Lemma~\ref{lem:trap} therefore applies with the following strip bounds:
\begin{center}\small
\begin{tabular}{ll}\toprule
Inner quantity & Valid $B_0$\\\midrule
$Z_f,Z_g$ & $2^{R(\sigma)},\ 2^{S(\sigma)}$\\
$Z_{f,\kappa},Z_{g,\kappa}$ & $2^{\kappa R(\sigma)},\ 2^{\kappa S(\sigma)}$\\
$\mathcal P(x_j)$ & $2^{R(\sigma)+S(\sigma)}$\\
$\mathcal Y(x_j)$ & $(R(\sigma)+S(\sigma))2^{R(\sigma)+S(\sigma)}$\\
$\langle2^{2r_\beta}\rangle$ &
 $2^{2\sum_j|r_j|\cosh(j\beta)\cosh(j\sigma)}$\\
$\langle2^{2s_\beta}\rangle$ &
 $2^{2\sum_j|s_j|\cosh(j\beta)\cosh(j\sigma)}$\\\bottomrule
\end{tabular}
\end{center}
Here $r_\beta(x)=\sum_jr_j\cosh(j\beta)\cos(jx)$; $s_\beta$ is
defined similarly. The node $x_j$ is real, so the same convolution bounds
hold uniformly for all outer nodes.
Outward interval arithmetic must enclose the exact node sums, including
the trigonometric arguments, exponentials, logarithms and all normalizers.
Enlarge each inner interval by its own $\mathcal E_N$ before performing
division or taking logarithms. Every denominator/logarithm interval used
must have a verified strictly positive lower endpoint. A larger precision
or agreement of two meshes is not a substitute for these error bounds.

\paragraph{The finite DFT acceleration is an exact sum identity.}
For node arrays $a_j=F(2\pi j/N)$ and $b_j=G(2\pi j/N)$ define
$\widehat a_\ell=\sum_j a_j e^{-2\pi i\ell j/N}$ and use the inverse
transform with prefactor $1/N$. Orthogonality of the finite geometric
sum gives, with all indices read modulo $N$,
\[
 \operatorname{DFT}^{-1}
       (\widehat a_\ell\widehat b_{-\ell})_j
 =\sum_t a_{j+t}b_t.
\]
Indeed, expansion of the left side retains exactly the terms with
$-j_a+j_b+j=0\pmod N$. Thus the raw convolution node average is this
inverse transform divided by $N$ once more. For $\mathcal Y$, add the
correlations of $(Fr,G)$ and $(F,Gs)$ and make the same division.
The identity applies to exact node values. Rigorous ball DFT arithmetic
must enclose that finite calculation; the analytic remainders above
still have to be added afterward. An ordinary floating FFT without
outward error control would not be a validated replacement.

\begin{lemma}[A genuine zero-free strip]\label{lem:strip}
Let $\beta>0$, and suppose rigorously computed bounds enclose
\[
 A_f=\frac{\langle2^{2r_\beta}\rangle}{Z_f^2},\quad
 A_g=\frac{\langle2^{2s_\beta}\rangle}{Z_g^2},\quad
 q_0=\sqrt{(A_f-1)(A_g-1)}<1.
\]
Then on $|\operatorname{Im}z|\le2\beta$,
\[
 |Q(z)-1|\le q_0,
 \qquad |L(z)|\le-\log_2(1-q_0),
\]
and $Q$ has a holomorphic logarithm real on the real axis. Moreover
\begin{equation}\label{eq:B-strip-bound}
 |B(z)|\le
 \frac{\sqrt{A_fA_g}(R(\beta)+S(\beta))}{1-q_0}
       +|\log_2(Z_fZ_g)|.
\end{equation}
\end{lemma}
\begin{proof}
The entire periodic functions $f,g$ are real and even on the real circle
because their defining trigonometric polynomials use only real cosines.
They have shifted means one by contour shifting. For a function real
on the real circle, its Fourier coefficients
satisfy $\widehat f(-k)=\overline{\widehat f(k)}$. Parseval therefore
writes the squared shifted $L^2$ norm as its constant term plus
nonnegative multiples of $\cosh(2ky)$, so it is nondecreasing in $|y|$.
This also bounds both positive and negative imaginary shifts.
At $y=\beta$ that squared norm is precisely $A_f$, since
$|F(t+i\beta)|^2=2^{2r_\beta(t)}$. In particular $A_f,A_g\ge1$.

For $z=x+iy$ with $|y|\le2\beta$, split the convolution contour:
\[
 Q(z)=\langle f(x+t+iy/2)g(t-iy/2)\rangle_t.
\]
Subtract the product of the shifted means and apply Cauchy--Schwarz
to $f-1,g-1$. Their squared norms are at most $A_f-1,A_g-1$;
Here the centered identity is
$\langle|f(\cdot+iy)-1|^2\rangle
=\langle|f(\cdot+iy)|^2\rangle-2\operatorname{Re}\langle f(\cdot+iy)\rangle+1
=\langle|f(\cdot+iy)|^2\rangle-1$, and similarly for $g$.
This proves the disk bound. The power series for $\operatorname{Log}(1+h)$,
with $h=Q-1$, supplies the holomorphic branch and gives
$|\operatorname{Log}Q|\le\sum_{j\ge1}q_0^j/j=-\ln(1-q_0)$.
Divide by $\ln2$ for the stated base-two bound. The strict inequality
$q_0<1$ also supplies a neighborhood of the closed strip.

Apply the same split contour to $\mathcal Y/(Z_fZ_g)$.
The absolute value of $r+s$ there is at most $R(\beta)+S(\beta)$;
Cauchy--Schwarz bounds the integral of $|fg|$ by $\sqrt{A_fA_g}$.
Since $|Q|\ge1-q_0$ and
$B=\mathcal Y/(Z_fZ_gQ)-\log_2(Z_fZ_g)$, this gives
\eqref{eq:B-strip-bound}. Real positivity alone would not justify
the complex logarithm used in this proof.
\end{proof}

\begin{proposition}[Per-certificate validated upper bound]\label{prop:validated}
For every filter occurring in a certificate or one of its comparison rows,
suppose Lemma~\ref{lem:strip}'s bounds are enclosed independently. Let
$L_v^*,B_v^*,H_v^*$ bound $|L_v|,|B_v|,|H_v|$ on the strip, and put
\[
 U_v^*=L_v^*+|\kappa_v-1|B_v^*+H_v^*,\qquad
 W_v^*=|2-\kappa_v|B_v^*+2L_v^*+H_v^*.
\]
For an old row vector $G_{v,i}$ form the following real-part upper bounds:
\begin{align*}
 \rho_0:&\quad (\Delta_b+U_v^*,\ \Gamma_b+W_v^*,\
                       (U_v^*+U_j^*)_{j\in C}),\\
 \rho_{1:8}:&\quad (A_bW_v^*+U_v^*,\ W_v^*-\Gamma_bL,\
                  \Gamma_bH+W_v^*,\ (U_v^*+U_j^*)_{j\in C}),\\
 \rho_9:&\quad (A_bW_v^*+U_v^*,\ W_v^*-15\Gamma_b,\
                       (U_v^*+U_j^*)_{j\in C}).
\end{align*}
Call the resulting vector $S_{v,i}$ and set
\[
 B_{v,i}^*=(2+2\cosh(2\beta))\,2^{z_{v,i}\cdot S_{v,i}}.
\]
If the interval $T_{v,i}$ encloses the average of the true outer node
values after all inner/arithmetic errors, then
\[
 I_{v,i}\le \sup T_{v,i}+\mathcal E_N(B_{v,i}^*,2\beta).
\]
Certification requires this upper bound to be strictly below an
independently enclosed lower bound for $2^{\eta_b-m_i}$.
\end{proposition}
\begin{proof}
Each affine row's real part is bounded by the displayed entry;
$A_b>0$ and all multipliers are nonnegative. Hence the real part of
the exponent is at most $z\cdot S$. The holomorphic extension uses
$2^z=\exp((\ln2)z)$ and the log branch already justified. Also
$|2+2\cos z|\le2+2\cosh(2\beta)$. Lemma~\ref{lem:trap}, followed
by the enclosing node interval, gives the result. Negative entries
such as $W_v^*-\Gamma_bL$ may be retained: they are valid upper bounds,
not assumed nonnegative rows at complex points.
\end{proof}

\paragraph{Exact target and pilot.}
The targets needed by the released theorem are $m_0=m_1=.00025$,
$m_2=\cdots=m_8=4$, and $m_9=17$, with $\eta_b=.51613$.
Kapila's common Proposition-B threshold alone is insufficient.
The bounded pilot consists of $(3,0),(3,1),(4,0)$ and at most one tail
certificate, using the actual source comparisons for each. A separate
record must classify every attempted pair as certified, unresolved or
refuted. A failed sufficient bound is unresolved, not a refutation;
a certified lower bound above threshold would be a refutation.
If bounds are rigorous and the measured cost permits the finite full
batch, the same contract can be applied to all 430 distinct pairs.

\section{Finite-circle convergence and effective evaluation}\label{sec:finite}
\begin{proposition}[All normalizers are part of the limit]\label{prop:finite}
Replace every inner average above by the normalized $P$-point average,
including the tilted normalizers. At the grid points the resulting
scores are precisely the finite-circle scores, and they converge
uniformly to the continuous scores. The corresponding outer finite
moments converge to their continuous moments. All convergence bounds
can be made effective for this fixed finite coefficient family.
\end{proposition}
\begin{proof}
For any continuously differentiable periodic real function $h$, comparison
on each interval of length $2\pi/P$ with its left endpoint gives
\[
 |T_Ph-\langle h\rangle|\le 2\pi\|h'\|_\infty/P.
\]
The same estimate holds uniformly in an additional real parameter
when its derivative bound is uniform. In this coefficient family,
$|r'|,|s'|\le2\sum_{j=1}^8j=72$. Thus valid bounds on the
integration-variable derivatives of $FG$ and $FG(r+s)$ are respectively
\[
 144(\ln2)2^{32},\qquad 144(32\ln2+1)2^{32}.
\]
For $F^\kappa,G^\kappa$ a bound is $144(\ln2)2^{32}$.
The raw normalizers and convolutions therefore differ from their
continuous versions by explicitly bounded constants divided by $P$.
Their positive lower bounds, for example $\mathcal P_P,\mathcal P\ge2^{-32}$
and $Z_{f,P},Z_f\ge2^{-16}$, hold on either circle without a limiting
argument. On these fixed positive compact intervals the quotient and
logarithm operations are Lipschitz; for the latter use
$|\log_2x-\log_2y|\le |x-y|/(\min(x,y)\ln2)$.
Substitution in~\eqref{eq:score-identities} gives computable constants
$C_u,C_w$ such that the uniform score errors are at most $C_u/P,C_w/P$.

For a fixed certificate, the exponent error is at most a computable
$C_E/P$. Its real exponents are uniformly bounded by another computable
constant $J$. The mean-value bound for $2^x$ on this interval, and
$0\le2+2\cos x\le4$, bound the change in the node integrand by a
constant divided by $P$. The continuous integrand is itself continuously
differentiable with a computable derivative bound, so applying the
Riemann estimate once more proves an explicit $C_I/P$ moment error.
All constants are finite expressions in the fixed rational coefficients,
$\pi,\ln2$ and elementary-function bounds, not input-dependent advice.
Consequently a certified continuous upper bound separated from a desired
finite threshold by $d>0$ gives the effective choice $P>C_I/d$.
There are finitely many modified certificates, so the maximum cutoff
works for all of them, with any strictly smaller finite margin.

For a prescribed fixed accuracy, choose a rational mesh from these
bounds and evaluate the continuous elementary functions by rational
Taylor series with explicit remainders. Bounded positive logarithm
arguments admit the convergent series for
$\ln y=2\sum_{j\ge0}((y-1)/(y+1))^{2j+1}/(2j+1)$.
Below the fixed cutoff compute finite averages directly; above it use
the continuous approximation and the uniform error bound. Approximating
the input residue angle uses polynomially many operations in $\log P$;
all meshes and accuracy constants are fixed. This gives an effective
uniform evaluation procedure, though not a claim that our Lean project
has already formalized this analytic argument.
\end{proof}

These arguments reconstruct the analytic steps in the pinned OpenAI
filtering/numerical sections and Kapila's finite-family adaptation
\cite{openai,kapila}. The proofs above supply their applicability to the
actual 43-filter definitions; the numerical moment thresholds remain
separate obligations until validated bounds are produced.

\section{All 430 original moments: validated numerical results}\label{sec:validated-results}
The full batch now supplies a validated enclosure for every exact pair
$(v,i)\in\{0,\ldots,42\}\times\{0,\ldots,9\}$, with no duplicate,
missing or inferred pair. All 430 enclosures satisfy
\begin{equation}\label{eq:validated-old-moments}
 I_{v,i}<2^{.51613-m_i},\qquad
 m_0=m_1=.00025,\quad m_{2:8}=4,\quad m_9=17.
\end{equation}
The accepted artifact is \texttt{evidence/all-moments.json}, SHA256
\begin{quote}\small
\nolinkurl{9646407eba82509f62c6464293d9fb472370d2e2c08a45ae07636c7bd0c84476}.
\end{quote}
Its source hash is
\nolinkurl{b397531cc5997e88c439d5290774861b46a34aa32d9cd71fa7dcced130359298},
the unchanged pinned Kapila input. Each row preserves its node average,
analytic outer error, enlarged integral enclosure, exact-threshold
enclosure, strict gap and the needed inner/strip conditions.

The implementation uses Python 3.12.14 and \texttt{python-flint 0.8.0},
160-bit ball arithmetic and one thread, with common half-width
$\beta=.08$ and inner width $\sigma=.1$.
The four-row pilot covered $(3,0),(3,1),(4,0),(4,9)$. The full run
used at most 8192 outer nodes and completed in 36.81 seconds, within
its 600-second and 131072-node ceilings. These are validation-batch
times, not subset-sum solver times. The numerical code was frozen in
implementation commits \texttt{398e7aa}/\texttt{f36b062}; the latter
also preserves direct-sum checks of the DFT convention.

For scale, the following deliberately outward rounded bounds summarize
the pilot pairs as present in the full artifact. They are not a
substitute for its other 426 separate certificates.
\begin{center}\small
\begin{tabular}{rrrr}\toprule
Pair $(v,i)$ & Nodes & Integral upper bound & Threshold lower bound\\\midrule
$(3,0)$ & 4096 & $1.385921$ & $1.429866$\\
$(3,1)$ & 512 & $1.426917$ & $1.429866$\\
$(4,0)$ & 1024 & $1.425608$ & $1.429866$\\
$(4,9)$ & 2048 & $3\cdot10^{-37}$ & $1.09109\cdot10^{-5}$\\\bottomrule
\end{tabular}
\end{center}

\paragraph{Independent replay and trust boundary.}
Procopius independently replayed every pair in a separate Python
environment in 36.14 seconds. Every numerical enclosure, threshold,
gap and node count matched; timing metadata naturally differed.
The review also checked exact Cartesian pair coverage, coefficient
and row correspondence, nonnegative multipliers, both tilted
normalizers, positive inner denominators, all comparison filters on
the common zero-free strip, and the explicit inner/outer remainders.
Direct finite-sum tests checked the library DFT sign and normalization;
separate missing, duplicate and altered source-record copies were
rejected. The numerical review is recorded in review commit
\texttt{7f1e398}, receipt \texttt{47a33c19}.

These are rigorous numerical results conditional on the correctness
of the source parser, Python bindings, FLINT/Arb ball elementary
functions and DFT implementation, together with the written analytic
error theorem proved here. They are not proofs checked inside Lean.
The enclosures include discretization and rounding errors; neither
mesh agreement nor a heuristic integration error estimate is their
justification. Internal agent review is not independent human peer
review or a site-awarded status.

The portable commands and exact environment are in
\texttt{moments/README.md}. From the repository root, the documented
pinned fetch places the input under \texttt{verification/evidence/};
\texttt{verification/moments/validate.py} with \texttt{--all} regenerates
all certificates. Use a new output filename when replaying so the
accepted artifact remains unchanged. Certificate evaluation is
independent of parameter discovery; no optimization is rerun.

\paragraph{Assumption change.}
Before this validation, all 430 inequalities in
\eqref{eq:validated-old-moments} were separate inherited assumptions.
They may now be replaced by the accepted validated evidence under
the computational trust just stated. The precise thresholds needed
by the released theorem are covered, including the stronger middle
and tail margins. No alias-coverage, checksum-estimator or complete
implementation theorem follows merely from those integral bounds.

\section{Reconstructing the changed certificate}\label{sec:boundary}
This proof starts from the two exact source records, not from the
released note's conclusion. Use zero-based filter indices and the
actual functions of Section~\ref{sec:analytic}. Kapila's Appendix~C gives
\begin{align*}
 E_{3,0}^{b}(x)={}&17.211(\Delta_b-u_3(x))
  +27.1977(\Gamma_b-w_3(x))\\
 &+8.3302(u_3(x)-u_5(x))+3.0053(u_3(x)-u_{16}(x)),\\
 E_{3,1}^{b}(x)={}&7.7547(A_bw_3(x)-u_3(x))
  +0(w_3(x)-\Gamma_b)+0(1.25\Gamma_b-w_3(x))\\
 &+4.2699(u_3(x)-u_5(x))+2.2635(u_3(x)-u_8(x))\\
 &+.461(u_3(x)-u_{16}(x))+3.4032(u_3(x)-u_{35}(x)).
\end{align*}
Here $\Delta_b=.02844$, $\Gamma_b=.45925$, $A_b=\Delta_b/\Gamma_b$.
In particular both zero entries are exact rationals; they are not small
floating-point estimates. The complete finite-data authentication must
identify exactly the expected pairs $(v,i)\in\{0,\ldots,42\}\times
\{0,\ldots,9\}$ and these comparison orders, not only a row count.

For the fixed published tuple put
\[
 \Delta=\eta-\tau,\quad\Gamma=2\tau-\eta,\quad
 A=\Delta/\Gamma,\quad q=\Gamma/\Gamma_b,\quad
 K=\Delta/q-\Delta_b.
\]
The numerical obligations include $\Delta>0$, $0<\Gamma<\Gamma_b$,
$K>0$, so $A-A_b=K/\Gamma_b>0$.
For pattern values $U_j,W_j$ select filter~3 only when $U_3\ge U_j$
for every $j$. Set $U=U_3,W=W_3$ and
$\alpha=1$ for $W\le\Gamma$, $\alpha=\Gamma/W$ otherwise.
Only $\alpha U<\Delta$ is the sampled case.
Define $E_{3,0}'$ by replacing the first two constants in $E_{3,0}^b$
by $\Delta,\Gamma$. Define $E_{3,1}'$ by replacing $A_b$ by $A$
and replacing the zero upper-row multiplier by $7.7547(A-A_b)$;
the lower-row multiplier remains exactly zero.

\begin{theorem}[Concrete domination and integral identities]\label{thm:boundary}
On sampled selected-filter-3 patterns, use $E_{3,0}'$ for
$W\le\Gamma$ and $E_{3,1}'$ for $\Gamma<W\le1.25\Gamma_b$.
Then, with exponents evaluated at the pattern values,
\[
 -\alpha U\le-\Delta+E_{3,i}'{}^{\pi}.
\]
For the actual analytic functions and the full normalized circle,
\begin{align*}
 \langle(2+2\cos x)2^{E_{3,0}'(x)}\rangle
   &=2^{\Lambda_0}\langle(2+2\cos x)2^{E_{3,0}^{b}(x)}\rangle,\\
 \langle(2+2\cos x)2^{E_{3,1}'(x)}\rangle
   &=2^{\Lambda_1}\langle(2+2\cos x)2^{E_{3,1}^{b}(x)}\rangle,
\end{align*}
where
\[
 \Lambda_0=17.211(\Delta-\Delta_b)+27.1977(\Gamma-\Gamma_b),
 \qquad \Lambda_1=7.7547(1.25)K.
\]
\end{theorem}
\begin{proof}
In the low case $\alpha=1$, $\Delta-U>0$, $\Gamma-W\ge0$,
and all listed comparisons are nonnegative. Since $17.211\ge1$,
$E_{3,0}'{}^\pi\ge\Delta-U$.
For the extended first band, $W>\Gamma>0$ implies $0<\alpha<1$,
$\alpha AW=\Delta$, and $AW-U>0$. The new upper multiplier is
nonnegative and its row $1.25\Gamma_b-W$ is nonnegative. The comparison
rows are nonnegative. Therefore
\[
 E_{3,1}'{}^\pi\ge7.7547(AW-U)\ge\alpha(AW-U)=\Delta-\alpha U.
\]
The lower row $W-\Gamma_b$ may be negative. Its contribution is zero
because its multiplier is exactly zero; no lower-boundary assumption
was silently added to this proof.

For every real $x$, direct subtraction gives
$E_{3,0}'-E_{3,0}^b=\Lambda_0$. In the first band the changed terms are
\[
 7.7547(A-A_b)w_3(x)
 +7.7547(A-A_b)(1.25\Gamma_b-w_3(x))
 =7.7547(A-A_b)1.25\Gamma_b=\Lambda_1.
\]
These are identities at every point of the circle, including those
unrelated to any selected pattern. Lemma~\ref{lem:exist} supplies actual
continuity/integrability. The identity $2^{a+b}=2^a2^b$ and linearity
of the normalized integral permit pulling the constant factor out.
No integration domain was shortened or selected by a pattern condition.
\end{proof}

\begin{lemma}[Ordinary transfer and complete coverage]\label{lem:ordinary}
For every other low record use $\Delta/q-u_v$ in place of
$\Delta_b-u_v$. For every bounded old band
$\Gamma_bL<W\le\Gamma_bH$, replace $A_b$ by $A$ and add
$z_0(A-A_b)$ to its upper-row multiplier. For the tail make only the
change $A_b\mapsto A$. These give sampled-pattern domination, with
pointwise shifts $z_0K$, $z_0KH$, and $z_0(A-A_b)w_v(x)$ respectively.
For filter~3 use Theorem~\ref{thm:boundary} in its first two regions.
Together the regions are disjoint and exhaustive, with every finite
upper endpoint assigned to the preceding region.
\end{lemma}
\begin{proof}
For an ordinary low pattern $W\le\Gamma_b$, $q\le\alpha\le1$.
If $U\ge0$, sampling and $\alpha\ge q$ imply
$\alpha U\le U$ and $U<\Delta/q$; more directly in either sign case
$\alpha U<\Delta\le\alpha\Delta/q$, so $\Delta/q-U>0$.
Also $\alpha\Delta/q\ge\Delta$. Since $z_0\ge1\ge\alpha$,
$z_0(\Delta/q-U)\ge\Delta-\alpha U$. The remaining rows are
nonnegative with nonnegative weights. The exponent difference is $z_0K$.

In a bounded old band, $W>\Gamma_bL\ge\Gamma_b>\Gamma$ and
$\alpha<1/L\le z_0$. Again $\alpha AW=\Delta$ and $AW-U>0$.
The old lower and upper rows are nonnegative; the upper multiplier
only increases. Domination follows. The $w_v(x)$ terms cancel in the
exponent difference, leaving $z_0(A-A_b)\Gamma_bH=z_0KH$.
The tail proof is identical except for the missing upper row. Its
shift is bounded above by $256z_0(A-A_b)$ using Lemma~\ref{lem:exist}.

The ordered old endpoints are
$\Gamma_b,1.25\Gamma_b,1.6\Gamma_b,2\Gamma_b,2.6\Gamma_b,
3.5\Gamma_b,5\Gamma_b,8\Gamma_b,15\Gamma_b$.
Replacing only filter~3's first endpoint by $\Gamma$ preserves order
because $0<\Gamma<\Gamma_b<1.25\Gamma_b$.
An ordered list partitions the real line into the initial closed ray,
successive intervals open on the left and closed on the right, and a
final open ray. Thus $W=\Gamma$ belongs to the low case, while
$W=1.25\Gamma_b$ belongs to the extended first band; no gap or
overlap occurs. All other endpoints and record assignments are unchanged.
\end{proof}

\paragraph{Connection to the fixed certificate.}
For a nonnegative old weight vector of sum at most $Z=418.36$, the
bounded-band update has sum at most $Z(1+A-A_b)<419$ once the fixed
scalar inequality has been proved. Other updates do not increase the
sum. The source maxima required for the five moment families are
$10.6428$ for the 42 ordinary low rows and $8.4537$ for the first bands;
the two special filter-3 rows have the shifts just proved. Every row's
pattern-to-coordinate-average error is at most $6\epsilon$ when
$A<1/4$: a first band row has error at most $(4A+3)\epsilon<4\epsilon$,
a comparison has at most $6\epsilon$, and low/boundary rows have at
most $4\epsilon$. Constants contribute no error and the zero row
contributes zero. This yields the released allowance $6Z'\epsilon$.
The $18\epsilon$ estimator allowance remains an algorithmic dependency,
not something implied by this row-error calculation.

The identities above transfer each old full-circle moment with its own
record. The fixed-data margin theorem must connect every one of those
records to the corresponding proved transfer; a generic lemma and an
unconnected table check would not establish that conjunction.

\section{Accepted concrete Lean checkpoints}\label{sec:formal-results}
The completed checkpoint is implementation commit \texttt{7f9548e},
with target \nolinkurl{Verification.DataBridge} importing
\texttt{Concrete}, \texttt{Boundary} and \texttt{Regularity}.
It uses Lean 4.24.0 with the following dependency pins:
\begin{quote}
Wharton: \nolinkurl{bed8e18f96db5955c1424abf3005b4fed4bd8d77}\\
mathlib: \nolinkurl{f897ebcf72cd16f89ab4577d0c826cd14afaafc7}.
\end{quote}
The full lockfile retains the other dependency pins.
The accepted theorem-to-manuscript mapping is deliberately scoped:
\begin{enumerate}
\item \textbf{Pattern domination.}
\nolinkurl{Verification.Published.low_pattern_domination} and
\nolinkurl{Verification.Published.extended_pattern_domination}
establish Theorem~\ref{thm:boundary}'s two domination inequalities
for the actual filter-3 records and published budgets. Their explicit
hypotheses are the relevant interval bounds, the listed nonnegative
comparisons and the sampled-pattern inequality; the desired domination
is not assumed. The exact zero lower multiplier is proved and used.
\item \textbf{Actual normalized-circle moments.}
\nolinkurl{Verification.Published.concrete_low_moment_identity} and
\nolinkurl{Verification.Published.concrete_band_moment_identity}
prove the full-circle shift conclusions of Theorem~\ref{thm:boundary},
together with integrability of both actual integrands and strict
positivity of the old moment. The definitions are the original
entropy-score integrals, not arbitrary functions called moments.
The measure is Lebesgue measure restricted to $[0,2\pi]$, with
the prefactor $(2\pi)^{-1}$ in the average. The positive-support proof
is inside that full interval. Continuity and positive denominators
are proved in \texttt{Regularity}, rather than assumed to avoid a
totalized-integral or totalized-logarithm artifact.
\item \textbf{Data bridges.}
\nolinkurl{filter_transcription}, \nolinkurl{certificate_transcription},
\nolinkurl{exact_pair_coverage}, \nolinkurl{source_filter_dimensions},
\nolinkurl{rawFilter_index_in_bounds} and
\nolinkurl{comparison_zip_no_truncation} in the same namespace
check the exact ordered 43-filter/430-record data and their dimensions.
The other bridge theorems give exact list lengths, comparison uniqueness
and the ordinary-low cap excluding filter~3. The literal data equalities
are kernel checked; correspondence of those literals to the authenticated
external TeX tables was also independently parsed and compared.
\item \textbf{Limited interval result.}
\nolinkurl{Verification.Published.complete_region_coverage} proves the
disjunction of the ten modified filter-3 intervals, with the indicated
open/closed endpoints. At this checkpoint it does \emph{not} prove
uniqueness, disjointness, or a region assignment for every ordinary
filter. Those broader claims have the written proof in
Lemma~\ref{lem:ordinary}, but their subsequent formal certificate links
are not credited to this checkpoint.
\end{enumerate}

\paragraph{Build, axioms and separate checks.}
The target and audit built from source, and Procopius independently
replayed all four completed modules using the pinned built imports.
The separately rerun transitive axiom report contains only
\texttt{propext}, \texttt{Classical.choice} and \texttt{Quot.sound}.
No accepted proof uses \texttt{sorry}, native-evaluation trust shortcuts
or a custom mathematical axiom supplying missing conclusions.
Explicit mathematical hypotheses were inspected separately from the
axiom report. The source/statement/axiom receipt is under
\texttt{evidence/boundary-checkpoint/}; independent acceptance is
review commit \texttt{097e859}, receipt \texttt{30e28528}.

The stronger all-import \texttt{lean4checker --fresh} attempt exceeded
its five-minute cap and is \emph{unfinished}, not passed. Successful
module source replay using built dependencies does not replace that
stronger imported-environment recheck. A separate automated Comparator
statement check was not run; the pinned Lake provides no such command.
Actual definitions and statements received independent meaning review.
No official challenge statement, website check or status upgrade is
claimed. The exact replay commands, elapsed times, hashes and timeout
are preserved in the checkpoint receipt.

\paragraph{Accepted fixed-certificate extension.}
Implementation commit \texttt{605cf983} adds \texttt{LogScalar},
\texttt{Scalar}, \texttt{RowMargins}, \texttt{Regions} and
\texttt{AllRows} to target \texttt{Verification}.
The combined source build and \texttt{AuditScalar.lean} pass.
Procopius independently replayed all five added modules and checked
the statements and permitted axioms: review commit \texttt{bd3cf784},
receipt \texttt{e3f55627}. Its accepted correspondence is as follows. The three
\nolinkurl{entropy_a}, \nolinkurl{entropy_twoa} and
\nolinkurl{entropy_hma} statements concern actual real logarithms
and binary entropy, not merely inequalities between supplied interval
endpoints. \nolinkurl{all_alias_gaps} checks the twenty fixed
alias and auxiliary gaps. \nolinkurl{all_430_transferred_margins} and
\nolinkurl{all_430_weight_sums} link the actual records to the
transfer allowances and rounding budget.
\nolinkurl{full_region_partition} states existence and uniqueness of
region assignment for every filter, including the boundary conventions.
These broader results are distinct from the earlier filter-3
interval disjunction.

The analytic link is
\nolinkurl{all_430_conditional_moment_bounds}: for each imported record,
it transfers the actual normalized-circle old moment to the new
threshold. Its explicit premises include the old moment inequality
and $w_v(x)\le256$ for every real $x$. We discharge these respectively
with the accepted outside-Lean certificates in
Section~\ref{sec:validated-results} and the written score proof in
Section~\ref{sec:analytic}; neither is a Lean conclusion of this theorem.
The reduced-score identities used by the numerical validation also
retain their written proofs. Moreover, per-record analytic
exponent/moment transfer is distinct from assembling sampled-pattern
domination in every ordinary region: the latter remains written,
while the new filter-3 pattern inequalities are already formalized.

Finally, \nolinkurl{exact_error_and_display_gap} and
\nolinkurl{fixed_tuple_ordinary_O} concern the fixed work exponent
$0.48767228002$, its strict allowance below $0.4876723$, and absorption
of the stated sublinear and logarithmic overhead. They do not prove
that the full algorithm satisfies that resource premise.
The 430 moment thresholds remain outside-Lean validated results,
and the full algorithmic obligations in Section~\ref{sec:remaining}
remain open. This extension retains the unfinished fresh-replay and
separate statement-check limitations described above.

\section{Two inherited algorithmic obligations reconstructed}\label{sec:algorithmic}
The following local results have complete written proofs. They do not
settle the separate alias-estimator or finite-precision arguments.

\begin{theorem}[Transformation and isolated bin]\label{thm:isolation}
Let $a_i$ be positive integers and $t\ge0$. Choose a power of two
$N_0\in[12n,24n)$, independent uniform $u_i\in\{1,\ldots,N_0\}$,
and put $C=1+nN_0$, $c_i=Ca_i+u_i$, $t_s=Ct+s$ for
$0\le s\le nN_0$. For every choice of the tie weights,
\[
 \exists I\subseteq[n]:a(I)=t
 \quad\Longleftrightarrow\quad
 \exists s,I:c(I)=t_s.
\]
On a NO input every transformed bin is empty. On a YES input the
probability that at least one transformed bin has exactly one index-set
representation is at least $11/12$.
\end{theorem}
\begin{proof}
If $a(I)=t$, take $s=u(I)\in[0,nN_0]$. Conversely,
$c(I)=Ct+s$ implies $C(a(I)-t)=s-u(I)$, whose absolute value is at
most $nN_0<C$. Since $a(I)-t$ is an integer, it is zero.

For completeness, let $\mathcal F$ be any nonempty family of subsets
of $[n]$, here the original solution family. For each coordinate $i$,
fix all $u_j$, $j\ne i$. If both classes exist, let $m_0$ be the
least tie sum among members not containing $i$ and $m_1$ the least
tie sum excluding $u_i$ among members containing $i$. Equality of the
two class minima requires $u_i=m_0-m_1$, at most one of its $N_0$
possible values. If either class is empty there is no such event.
If the overall minimum is attained by two different subsets, some
coordinate distinguishes them, so that coordinate's class minima
are equal. A union bound, conditioning on the other weights for each
coordinate, gives probability at most $n/N_0\le1/12$ of a nonunique
minimum. A unique minimum solution has a unique tie sum within the
solution family, and hence a unique transformed bin by the equivalence
already proved. Nothing excludes repeated $a_i$ or the empty subset.
\end{proof}

\begin{lemma}[Deficiency on uniform half-blocks]\label{lem:deficiency}
For fixed transformed weights define
\[
 \Sigma_c(V)=\{c(I):I\subseteq V\},\qquad
 d_c(V)=|V|-\log_2|\Sigma_c(V)|.
\]
Inserting one index changes deficiency by a value in $[0,1]$.
Replacing one index by another changes it by at most one in absolute
value. For a uniformly sampled $k$-subset $V_k$,
\[
 \operatorname{Var}(d_c(V_k))\le k/4\le n/4.
\]
Also $f(k)=\mathbb E d_c(V_k)$ is nondecreasing in $k$.
\end{lemma}
\begin{proof}
For $i\notin V$,
$\Sigma_c(V\cup\{i\})=\Sigma_c(V)\cup(c_i+\Sigma_c(V))$.
Its cardinality lies between $|\Sigma_c(V)|$ and
$2|\Sigma_c(V)|$, proving the insertion claim. Two sets differing
by an exchange both arise by inserting an index into their common
part, so their deficiencies lie in an interval of length one.
Sampling a uniform $(k+1)$-subset and deleting a uniform member
couples the uniform sizes $k+1,k$; monotonicity follows.

Expose the first $k$ members $X_1,\ldots,X_k$ of a uniformly random
ordered sample without replacement, and let $M_j$ be the conditional
expectation of the final deficiency after $j$ exposures. At a fixed
prefix, compare next choices $a,b$. Transposing $a,b$ gives a bijection
between the possible remaining completions. The corresponding final
sets are either identical (both indices eventually occur) or differ
by one exchange. Hence the two conditional expectations differ by at
most one. Conditional on the prefix, $M_j$ therefore lies in an
interval of length at most one. A real random variable in an interval
of length one has variance at most $1/4$: expand
$(X-\min X)(\max X-X)\ge0$ and complete the square in its mean.
Thus each martingale difference has conditional second moment at most
$1/4$. Distinct differences are orthogonal by iterated conditional
expectation. Summing their variances gives the bound $k/4$.
\end{proof}

\begin{theorem}[A capped preliminary search and the LD split]\label{thm:preliminary}
Let $n\ge2$, $b\ge1$, and use the fixed published $0<\tau<1/2$.
Fix the tie weights and $\zeta>0$. The condition
\begin{equation}\label{eq:LD}
 \mathbb E d_c(V_{\lfloor n/2\rfloor})
       \le(\tfrac12-\tau+\zeta)n
       \tag{LD}
\end{equation}
is exactly the low-deficiency condition used here. There is a
preliminary procedure returning an exact decision or CONTINUE, using
$O(2^{\tau n}\operatorname{poly}(n+b+2))$ word operations on every
history. If LD fails, for all sufficiently large $n$ it returns
CONTINUE with probability at most
\[
 \frac{2}{\zeta^2 n}+n2^{-n^2}.
\]
This conclusion concerns a fixed-input probability; an observed
CONTINUE is not a certificate that LD holds.
\end{theorem}
\begin{proof}
Sample a uniform half-block $V$ and its complement, of sizes
$\lfloor n/2\rfloor$ and $\lceil n/2\rceil$. If LD fails, either
block $U$ has expected deficiency greater than
$(\tfrac12-\tau+\zeta)n$, by Lemma~\ref{lem:deficiency}.
Chebyshev's inequality with variance at most $n/4$ gives
\[
 \Pr\{d_c(U)<(\tfrac12-\tau+\zeta/2)n\}
 \le \frac1{\zeta^2n}.
\]
Outside these two bad events,
\[
 \log_2|\Sigma_c(U)|=|U|-d_c(U)
 \le\tau n-\zeta n/2+1/2<\tau n
\]
when $n>1/\zeta$. Both lists of distinct half sums fit within
$B_n=2^{\lceil\tau n\rceil}$.

Construct such a sorted list by starting with $\{0\}$ and, for each
weight, merging the current sorted list with its translate, deleting
duplicate values. If the next distinct output would exceed $B_n$,
return CONTINUE before writing it. The number of distinct sums is
nondecreasing under insertions, so if the final list fits, every prefix
fits. Each merge costs $O(B_n)$ operations. Once both full lists fit,
the sorted two-pointer search answers exactly whether some left and
right sums total any required $t_s$. There are $O(n^2)$ targets, so
even checking them separately costs only a polynomial times $B_n$.
If either list hits its cap, no inexact decision is emitted.

For an explicit bounded sampler, expose a uniform ordered sample
without replacement using partial Fisher--Yates swaps. A required
integer in $\{0,\ldots,m-1\}$ is obtained by masking a fresh uniform
word to $\lceil\log_2m\rceil$ bits and rejecting values at least $m$.
Acceptance is at least $1/2$. Cap each draw at $n^2$ proposals and
return CONTINUE on exhaustion. Couple this with the ideal unbounded
sampler: at most $n$ draws are made, so discrepancy probability is at
most $n2^{-n^2}$. The ideal subset is exactly uniform. The bounded
sampler has polynomial cost on every history.

All sums have $O(b+\log n)$ bits: $N_0<24n$ and $C<24n^2+1$
give $\sum_i c_i<31n^3 2^b$ for $n\ge2$, $b\ge1$.
The bit bound $b+3\log_2n+5$ is below the prescribed
$\lceil4(n+b+\log_2(n+2))\rceil$ word size on this domain, so sums
and keys occupy a fixed number of words. Two capped final lists plus
the alternating merge workspace can be kept in at most $4(B_n+1)$
words: the translated list is read implicitly, not copied. The partial
Fisher--Yates array, input and counters require at most $K n$ additional
words for a fixed implementation constant $K$. Thus simultaneous
storage is at most $4(B_n+1)+K n$, whose address length is
$\tau n+O(\log n)$. Since $\tau<1/2$, this is below the prescribed
word size for all $n$ above a fixed implementation cutoff. For the
finitely many smaller $n\ge2$, enumerate subsets with a running sum
in $O(n)$ storage and answer exactly; its $2^n$ time contributes only
a fixed factor to the stated $O(2^{\tau n}\operatorname{poly}(n+b+2))$
bound. This cutoff is independent of the input bit length and of $c$.
Fixed rational $\tau$ permits exact calculation of its ceiling.
Every allocation and loop has an explicit cap,
including the failure histories, proving the stated all-history cost.
\end{proof}

These proofs reconstruct the transformation and preliminary-search
results of the pinned OpenAI setup and Kapila Sections~2.1/4
\cite{openai,kapila}. They establish the local statements under their
displayed hypotheses without assuming an exact subset-sum solver or
an LD oracle. In the composed analysis one fixes the tie weights first,
splits according to LD, and draws the main procedure's independent
randomness whether or not an earlier exact answer makes it unused.
This avoids conditioning the main proof on the observed continuation.

\section{The consequential algorithmic dependencies that remain}\label{sec:remaining}
\begin{proposition}[Exact alias classes and sound insertion]\label{prop:alias-sound}
Use the actual integer key $(c(A),b(A))$ for each half-block subset,
where $b_i=c_i\bmod P$ is the least nonnegative residue. Keep a
dictionary key $(\kappa_L,\kappa_R,Z)$ only once, and store the value-map
phase sums $\Phi_L,\Phi_R$ over all subsets with the respective key.
Every insertion made by Kapila's Section~5.2 matching/expansion rule
is an alias class. If all alias classes are reached, the exact output
equals the full alias phase sum, with no multiplicity error from
repeated discoveries or repeated original weights.
\end{proposition}
\begin{proof}
For precision, $T_{s,q}=(t_s\bmod P)+qP$, $0\le q<n$.
Every modular representation has exactly one such carry: its integer
sum of reduced weights is in $[0,n(P-1)]\subset[0,nP)$ and is congruent
to the least residue of $t_s$. Subtracting that residue and dividing
by $P$ gives an integer in $\{0,\ldots,n-1\}$.
Write $c^F_i=(-1)^{[i\in F]}c_i$ on the large block, and similarly
for $b^F$. Directly partitioning $Z$ into $Z\setminus F$ and
$Z\cap F$ proves
$c^F(Z)=c(Z\triangle F)-c(F)$, and the same identity for $b$.
In an insertion the masks $X,Y$ are disjoint; put
$Z=(X\cup Y)\triangle F$. The integer matching condition is
\[
 b(\kappa_L)+b^F(X)+b(\kappa_R)+b^F(Y)
       =T_{s,q}-b(F),
\]
where $T_{s,q}\equiv t_s\pmod P$. The sign identity gives
$b(A\cup Z\cup B)=T_{s,q}$ for every half-block pair having those
keys, and thus its unreduced sum is congruent to $t_s$ modulo $P$.
The skipped exact interval in the matching rule excludes precisely
\[
 c(\kappa_L)+c^F(X)+c(\kappa_R)+c^F(Y)=t_s-c(F).
\]
The same identity therefore shows that each inserted member has
unreduced sum different from $t_s$. It is an alias, not an exact
representation.

Every index set $I$ belongs to exactly one class:
$(\kappa(I\cap L),\kappa(I\cap R),I\cap M)$, for the fixed disjoint
partition $L,M,R$. Every member of a class has the same unreduced
and reduced sums. Additivity of the retained phases gives
\[
 \sum_{I\text{ in class}}\chi(I)
 =\left(\sum_{A:\kappa(A)=\kappa_L}\chi(A)\right)
  \left(\sum_{B:\kappa(B)=\kappa_R}\chi(B)\right)\chi(Z)
 =\Phi_L(\kappa_L)\Phi_R(\kappa_R)\chi(Z).
\]
Distinct class keys partition the relevant index sets, and the
dictionary stores each key at most once. Consequently reaching every
alias class suffices for the claimed exact sum. This does not prove
that the randomized procedure reaches every class before its cap;
that coverage/cost/probability obligation remains below.
\end{proof}

The following are exact obligations, not conclusions of the local
proofs or of scalar certificate checking.
\begin{enumerate}
\item \textbf{Alias coverage and cost.} Under LD, with the published
block sizes and slack, Kapila's Section~5 procedure must return every
modular alias contribution in its deduplicated value classes with the
retained signing, except with a uniformly small failure probability.
Its expected-size estimates must imply the stated deterministic caps
through the actual abort policy, without losing the per-input success
bound. The half-block bad-prime bound uses
$\delta=(\eta-(\mu/2)\log_2 3)/2$; the actual requirements include
$\delta>1/2-\tau+\zeta$ and the bit-length guard slack $\theta<\delta$.
\item \textbf{Same-phase cancellation.} For each target and modular
prime, the alias sum and the checksum must use identical stored phases.
Every modular representation with a different integer sum must occur
exactly once in the deduplicated class sum. This gives exact residual
zero on NO inputs and a modulus-one residual at an isolated bin;
it does not by itself bound either numerical estimator's error.
\item \textbf{Block-table estimator.} For the actual finite-circle
density functions, every selected sampled pattern must satisfy the
table-coverage/second-moment hypotheses in OpenAI Section~4 as adapted
by Kapila Section~6. The released $18\epsilon$ allowance belongs here.
The distribution used to factor coordinate moments must be the original
independent phase law, not a law conditioned on a favorable pattern,
table event, alias-success event or CONTINUE.
\item \textbf{Finite precision and uniform constants.} Scores and
elementary functions must be evaluated with enough accuracy for the
actual block thresholds. Stored phases must remain consistent across
both subroutines; all accumulated errors must stay below the decision
threshold gap. Section~\ref{sec:finite} supplies effective fixed-score
evaluation, but not the full propagated complex-arithmetic error bound.
The program must generate required constants/tables, not receive them
as advice.
\item \textbf{Every-history resources.} Prime search, table construction,
trial lists and visits, rational arithmetic, random draws, memory and
indirect addresses all require deterministic caps with charged word
cost. A finite expected work estimate alone is insufficient. The
published scalar gaps verify proposed exponents, not that a described
procedure actually obeys each one. The guarded program and exact
fallback must be a single program independent of the polynomial-bit
parameter $c$.
\end{enumerate}

\paragraph{Precise cancellation fact.}
If $\chi(I)$ is the same retained unit-modulus phase in both sums,
then, identically,
\[
 \sum_{c(I)\equiv t_s\ (\mathrm{mod}\ P)}\chi(I)
 -\sum_{\substack{c(I)\equiv t_s\ (\mathrm{mod}\ P)\\c(I)\ne t_s}}\chi(I)
 =\sum_{c(I)=t_s}\chi(I).
\]
This is a finite partition identity, valid with repeated original
weights because subsets are index sets. It proves the cancellation
\emph{if} the two enumerated/sampled quantities correspond to those
exact sums with common phases. Completeness of the class algorithm and
accuracy of the estimates remain separate obligations; renaming their
outputs by these sums would assume the desired conclusion.

The current supplement therefore does not claim that the complete
word-RAM algorithm has been independently verified. Any unresolved
point above remains an explicit inherited dependency even after the
boundary, scalar and continuous-moment obligations are settled.

\subsection{A bounded precision obligation that can be discharged locally}
\begin{proposition}[Alias-sum arithmetic, given the correct dictionary]
\label{prop:alias-precision}
Suppose the phase maps and dictionary are those of
Proposition~\ref{prop:alias-sound}, with coverage already supplied.
Use the same exact retained phase vector modulo $P$ throughout.
Compute each unit phase to absolute complex error at most
$\delta_0=2^{-2n^2}$, and round each real and imaginary component
after addition or multiplication to absolute error at most
$R=2^{-n^2}$. Then the total alias-output error is at most
\[
 32\,2^{-n^2+3n}<1/16\qquad(n\ge8).
\]
These operations can be performed with polynomially many charged word
operations per phase or arithmetic operation. This is not a proof of
coverage or of the separate checksum estimator's precision.
\end{proposition}
\begin{proof}
There are at most $2^n$ summands in either half-map and at most $2^n$
dictionary keys: distinct dictionary classes are nonempty and disjoint
subsets of the $2^n$ possible index sets. The Euclidean error of a
complex rounding is at most $2R$. Thus either computed half-map sum
differs from its exact value by at most
$2^n\delta_0+2^{n+1}R\le3\,2^nR=:e$.
The exact sum has modulus at most $2^n$, and its approximation has
modulus at most $2^{n+1}$ for $n\ge2$.

In one class product $\Phi_L\Phi_R\chi(Z)$, replacing the factors
by their approximations changes the value by at most
\[
 e2^n+2^{n+1}e+2^{2n+2}\delta_0
 \le13\,2^{2n}R.
\]
Two rounded multiplications contribute at most $6R$, because the
approximate unit phase has modulus at most two. Hence a class term
has error at most $20\,2^{2n}R$. Summing at most $2^n$ terms and
rounding the additions gives error at most
$20\,2^{3n}R+2^{n+1}R\le32\,2^{3n}R$. The displayed strict
$1/16$ bound holds already at $n=8$ and improves thereafter.
All estimates are deterministic and uniform over the retained phases.

For an explicit effective phase evaluation, first sum the integer
residues modulo $P$ exactly, then approximate the bounded angle
$2\pi k/P$. The identity
$\pi/4=\arctan(1/2)+\arctan(1/3)$ and the alternating arctangent
series give a rational $\pi$ enclosure to error $2^{-8n^2}$ using
$16n^2$ terms. Rational division at the same precision gives an angle
error $O(2^{-8n^2})$. Evaluate sine and cosine by their Taylor
polynomials through degree $16n^2$, on $|x|\le8$. The remainder
is bounded by the exponential-series tail
\[
 \frac{8^{N+1}}{(N+1)!}\,
       \frac1{1-8/(N+2)},\qquad N=16n^2.
\]
For $n\ge8$, use $N!\ge(N/2)^{N/2}$ and $N\ge1024$;
this bound is at most $16\,2^{-24n^2}$.
The angle error and Taylor errors together are strictly less than
$2^{-2n^2}$ in complex norm, with ample slack.

All exact input phases have $O(n)$ bits because the published prime
interval has $\log_2P=O(n)$. Real and imaginary intermediate magnitudes
are at most $2^{O(n)}$, even when summing all dictionary terms.
Keeping $n^2$ fractional bits therefore uses polynomially many bits,
as do the rational Taylor calculations. Elementary multiword integer
arithmetic implements each operation with polynomially many word
operations, charging retained storage and rounding. No transcendental
oracle or uncharged table is required. Finite smaller $n$ can use
the exact fallback already described.
\end{proof}


\section{Fixed data, formalization and the allowable conclusion}
The published tuple, all 43 filters and the exact expected 430
filter/region pairs are fixed throughout this phase. A kernel proof of
rational comparisons must be distinguished from a proof that those
rationals enclose actual real logarithms or entropy. Likewise a generic
transfer theorem must be instantiated with the actual functions and
coefficient records; naming an arbitrary quantity a moment is insufficient.
The written identities in this supplement do not by themselves certify
an implementation, a Lean theorem, or an integration output.

The internal work exponent is exactly $\beta_{\rm work}=.48767228002$ and the
advertised exponent is $\alpha=.4876723$, with
$\alpha-\beta_{\rm work}=1.998\cdot10^{-8}$. If an algorithmic argument supplies
the actual all-history bound
$2^{\beta_{\rm work}n+O(n^{.7})}\operatorname{poly}(n+n^c)$ for each fixed $c$,
then it is $O_c(2^{\alpha n})$: taking base-two logarithms of the extra
factor gives $O(n^{.7})+O_c(\log n)=o(n)$, eventually below
$(\alpha-\beta_{\rm work})n$. This is an asymptotic implication, not a proof of
the premise's resource accounting or a license to omit that premise.

\begin{theorem}[The strengthened conditional implication]\label{thm:strengthened}
Retain the published fixed tuple and its exact scalar certificate.
Assume the remaining source algorithmic claims identified in
Section~\ref{sec:remaining}: the averaged-layer/alias coverage and
cap analysis of Kapila Sections~3 and~5, the block-table/checksum
probability analysis of OpenAI Section~4 as adapted in Kapila Section~6,
and their finite-precision/full-program implementation claims with the
published parameters, guard and slack. Under the explicit validated
arithmetic trust of Section~\ref{sec:validated-results}, the same single
randomized word-RAM program halts on every input and is correct with
probability at least $2/3$ on every input. For every fixed positive
integer $c$ its time on every random execution with $b\le n^c$ is
$O_c(2^{.4876723n})$.

The 430 original continuous moment inequalities are no longer additional
unverified premises: they are supplied by the validated certificates.
This remains a conditional algorithmic theorem, not a full Lean proof
or an independently verified complete subset-sum algorithm.
\end{theorem}
\begin{proof}
Equation~\eqref{eq:validated-old-moments} gives exactly the old margins
needed by the five released transfer families. The concrete transfer
proofs here and the unchanged exact scalar certificate supply a finite
margin $10^{-7}$; the prescribed rounding loss is bounded by
$2532\cdot10^{-12}$, leaving $9.7468\cdot10^{-8}>0$.
Section~\ref{sec:finite} supplies the needed effective finite-circle
approximation. These statements do not invoke an estimator success
conclusion as a scalar or analytic hypothesis.

For fixed tie weights, split according to LD, not according to an
observed CONTINUE. In the complementary case the preliminary procedure
answers exactly except with probability $O(1/n)$; under LD the
remaining alias and checksum hypotheses give the required uniform
accuracy with failure probability $o(1)$. Use the same retained phases
in both quantities. Sound class summation and exact cancellation leave
zero on a NO input. On a YES input the isolation proof gives a unique
exact bin with probability at least $11/12$, whose phase has modulus
one. A combined approximation error less than $1/4$ separates the two
cases at modulus threshold $1/2$. The total failure is at most
$1/12+o(1)$, hence below $1/3$ above a fixed cutoff; finite exceptions
use exact enumeration. The separate alias-arithmetic bound proves its
local part of this accuracy claim, but the checksum precision and
coverage parts are explicitly retained source hypotheses.

The stated implementation hypotheses concern the specified subroutines
and their deterministic caps, not a renamed assertion of the final
decision theorem. They give the all-history bound with exponent
$\beta_{\rm work}$ above. The input guard
$10^{12}\lceil\log_2(b+1)\rceil\le n$ holds eventually for every
fixed $b\le n^c$; a guard failure invokes the exact fallback.
For each fixed $c$ those exceptions occur at only finitely many $n$,
and their cost is absorbed into the constant $C_c$. The program and
its cutoff/guard do not depend on $c$. Strict exponent absorption then
gives the displayed ordinary-$O$ bound. Every-history cost and eventual
probability are used separately; an expected-time bound is insufficient.
\end{proof}

\paragraph{What has and has not been removed.}
Compared with the published assumptions, actual score regularity/bounds,
effective finite-circle approximation, isolation and the preliminary
LD split now have reconstructed reviewed written proofs. Sound alias
insertion, class summation and its local arithmetic error have written
proofs as well. The entire old-moment family is validated outside Lean.
The unproved averaged-layer/alias coverage and cap chain, checksum
estimator/precision chain and complete implementation/resource chain
remain as precisely listed source hypotheses. The new fixed-certificate
Lean development is assessed separately; this numerical result alone
does not settle its formal obligations or imply any site status.

Only exact obligations supported by completed evidence may be removed
from the proposed strengthened conditional statement. In particular,
unresolved analytic hypotheses in a formal theorem remain hypotheses
even when its axiom list is clean. The allowed ordinary Lean foundational
axioms are reported separately from explicit mathematical premises.
Unfinished or experimental modules cannot enter the accepted dependency
chain by assertion, a custom axiom or a native-evaluation trust shortcut.

\paragraph{Provenance and assistance.}
The density and numerical-convolution framework is OpenAI's; the 43-filter
and alias construction is Annanay Kapila's; the earlier transfer framework
and finite-data formalization are Joshua Wharton's
\cite{openai,kapila,wharton}. The present supplement reconstructs the
specific applicability and proofs stated here. The intended author is
Yanni Bills, assisted by Researcher-Ed\_Lorenz for source/proof analysis,
Computational-Richardson for formal and validated numerical work,
Reviewer-Procopius for adversarial internal review, and the Coordinator
for integration. Internal agent review is not independent human peer
review or official website verification. Nothing in this supplement
authorizes publication, a release change or a verification-status request.

\bibliographystyle{plain}
\bibliography{verification-references}
\end{document}
